\section{The local terms and torsion in the fixed-base character reduction}
\label{sec:fixed-base-local-terms}

Use \(K,H,G,D,N,O_\ell,b_\ell\) from AC-K25.
Let \(\mathcal V\) be the finite places of \(K\) ramified
in \(H/K\), and put
\[
 r=\#\mathcal V,\qquad
 \mathcal R=\sum_{v\in\mathcal V}\log Nv,\qquad
 c_0=\log2+\tfrac12\log3 .
 \tag{FL.1}
\]
For these hosts \(r\leq12\).
For \(e\geq1\) and \(\eta\in\widehat G\), write
\[
 A_{\eta,e}=(O_\ell\otimes_{\mathbb Z_\ell}
                  \mathbb Q_\ell/\mathbb Z_\ell)(\eta)(-e),
 \qquad
 T_{\eta,e}=O_\ell(\eta^{-1})(e+1).
 \tag{FL.2}
\]
The precise Tate dual lattice is initially
\(\operatorname{Hom}_{\mathbb Z_\ell}(O_\ell,\mathbb Z_\ell)
(\eta^{-1})(e+1)\).
The codifferent is a free rank-one \(O_\ell\)-module.
Choosing one of its generators identifies this lattice with
\(T_{\eta,e}\), compatibly with its Galois action.
Thus all finite group orders below use the actual dual
coefficient.  Dividing their logarithms by \(b_\ell\),
rather than by the residue degree alone, retains the
ramification in the coefficient ring.

\begin{theorem}[AC-K26-local: a polynomial bound for all local comparison terms]
\label{AC-K26-local}
Let \(\mathcal S_{\ell,\eta}(e)\) be the strict group
in (FC.2), and set
\[
 \mathcal H_{\ell,\eta}(e)
             =H^2(G_{K,\Sigma_\ell},T_{\eta,e}).
 \tag{FL.3}
\]
These \(H^2\) groups are finite.
For each fixed \(H,e\), the following sums are finite,
and one has the explicit bound
\[
 \begin{split}
 &\sum_\ell\frac1{b_\ell}\sum_{\eta\in\widehat G}
  \left|\log\#\mathcal H_{\ell,\eta}(e)
                  -\log\#\mathcal S_{\ell,\eta}(e)\right|\\
 &\hspace{8mm}\leq
      16eD^2\log(8eD)+De\mathcal R+2Dr c_0 .
 \end{split}
 \tag{FL.4}
\]
The statement is unconditional.
In particular the difference between the two sums
without absolute values obeys the same bound.
On the actual hosts and weights of (FG.3), summing
these errors over \(1\leq e<t_x=\lfloor x^{401}\rfloor\)
gives
\[
 O_\varepsilon\!\left(
                  x^{806+16\varepsilon}(\log x)^2\right)
          =o\!\left(\exp(x^\varepsilon/2)\right).
 \tag{FL.5}
\]
\end{theorem}
\begin{proof}
\emph{1. An elementary cyclotomic invariant estimate.}
Let \(U\) be an open subgroup of
\(\mathbb Z_\ell^\times\) of index \(I\).
For any positive integer \(w\), let
\(B_w=(\mathbb Q_\ell/\mathbb Z_\ell)(w)^U\).
The same group order is obtained for twist \(-w\).
It is finite; if it is nonzero and has order \(\ell^k\),
then
\[
                           \ell^k\leq4wI .
 \tag{FL.6}
\]
For odd \(\ell\), write
\[
 U=\mu_m\times(1+\ell^s\mathbb Z_\ell),\qquad
 m\mid(\ell-1),\quad s\geq1,\quad
 I=\frac{\ell-1}{m}\ell^{s-1}.
 \tag{FL.7}
\]
The product description follows from the prime-to-\(\ell\)
and pro-\(\ell\) factors of \(\mathbb Z_\ell^\times\):
a closed subgroup splits into those factors by taking
limits of suitable powers.
If \(m\nmid w\), a torsion unit gives a unit value of
\(u^w-1\), so \(B_w=0\).
If \(m\mid w\), the binomial lifting-of-exponent formula
on \(1+\ell^s\mathbb Z_\ell\) gives
\(k=s+v_\ell(w)\).
Consequently
\[
 \ell^k
    =\frac{\ell}{\ell-1}I\,m\,\ell^{v_\ell(w)}
    \leq\tfrac32 Iw ,
 \]
because \(m\ell^{v_\ell(w)}\) divides \(w\).
For \(\ell=2\), write
\(U\cap(1+4\mathbb Z_2)=1+2^s\mathbb Z_2\), \(s\geq2\).
Its index in \(1+4\mathbb Z_2\) is \(2^{s-2}\leq I\).
The element \(1+2^s\) belongs to \(U\) and
\[
              v_2((1+2^s)^w-1)=s+v_2(w).
 \]
It follows that the invariant group is killed by
\(2^{s+v_2(w)}\leq4Iw\).
In all cases the invariants form a subgroup of the cyclic
group \(\ell^{-k}\mathbb Z_\ell/\mathbb Z_\ell\), which
proves finiteness and (FL.6).
No splitting of a general \(2\)-adic subgroup into a sign
factor and a principal-unit factor has been assumed.

If \(M\) is a finite extension of \(\mathbb Q_\ell\), its
cyclotomic image has index at most \([M:\mathbb Q_\ell]\).
The analogous global index for a number field is at most
its degree.  In either case this follows by applying the
surjective cyclotomic character of the corresponding
rational absolute Galois group to a subgroup of finite
index: the index of the image cannot increase.
This also proves these bounds when the field intersects
the cyclotomic extension nontrivially.

\emph{2. Full localization versus the inertia condition.}
Define
\[
 \operatorname{Sha}^1(A_{\eta,e})=
 \ker\!\left(H^1(G_{K,\Sigma_\ell},A_{\eta,e})
           \longrightarrow
              \bigoplus_{v\in\Sigma_\ell}
                            H^1(K_v,A_{\eta,e})\right).
 \tag{FL.8}
\]
Modified cohomology at the complex place is zero.
This full localization kernel is a subgroup of
\(\mathcal S_{\ell,\eta}(e)\), which is finite by AC-K25.
Inflation--restriction at a finite \(v\nmid\ell\) identifies
\[
 H^1_{\rm unr}(K_v,A_{\eta,e})
   :=\ker\!\left(H^1(K_v,A_{\eta,e})
                        \longrightarrow H^1(I_v,A_{\eta,e})\right)
   = A_{\eta,e}^{I_v}/(\mathrm{Fr}_v-1)A_{\eta,e}^{I_v}.
 \tag{FL.9}
\]
Here \(\mathrm{Fr}_v\) is arithmetic Frobenius.
For a procyclic residue-field Galois group the displayed
degree-one cohomology formula holds for these discrete
\(\ell\)-primary coefficients.
Thus localization induces an injection
\[
 \mathcal S_{\ell,\eta}(e)/\operatorname{Sha}^1(A_{\eta,e})
       \lhook\joinrel\longrightarrow
        \bigoplus_{v\in\mathcal V,\ v\nmid\ell}
                             H^1_{\rm unr}(K_v,A_{\eta,e}).
 \tag{FL.10}
\]
At \(v\mid\ell\) both kernels require full vanishing.
All other finite places in \(\Sigma_\ell\) belong to
\(\mathcal V\).

Compact Poitou--Tate duality and local Tate duality give
\[
 \operatorname{Sha}^2(T_{\eta,e})\simeq\operatorname{Sha}^1(A_{\eta,e})^\vee,\qquad
 H^2(K_v,T_{\eta,e})\simeq H^0(K_v,A_{\eta,e})^\vee .
 \tag{FL.11}
\]
The exact source is Sharifi, \emph{Iwasawa Theory},
Appendix A, Theorems A.0.3 and A.0.6, with the
compact-coefficient extensions in Remarks A.0.4 and A.0.9;
the relevant localization sequence is Theorem A.0.8.
These are the same duality statements used in AC-K10,
now with the larger finite set \(\Sigma_\ell\).

Once the local invariant groups in (FL.11) are finite,
the localization map in degree two shows that
\(\mathcal H_{\ell,\eta}(e)\) is finite: its kernel is
the finite group in (FL.11), and its image is a subgroup
of the finite sum of local \(H^2\)'s.
Taking orders in the two localization kernels then gives
\[
 \begin{split}
 &\left|\log\#\mathcal H_{\ell,\eta}(e)
                  -\log\#\mathcal S_{\ell,\eta}(e)\right|\\
 &\quad\leq
       \sum_{v\in\Sigma_{\ell,\rm fin}}
                              \log\#H^0(K_v,A_{\eta,e})
       +\sum_{v\in\mathcal V,\ v\nmid\ell}
                              \log\#H^1_{\rm unr}(K_v,A_{\eta,e}).
 \end{split}
 \tag{FL.12}
\]
We prove the needed finiteness and sum these terms next.

\emph{3. The places over the coefficient prime.}
Fix \(v\mid\ell\), and choose a place \(w\) of \(H\)
above \(v\).  The character \(\eta\) is trivial on
\(G_{H_w}\), so restriction of invariants is an injection
\[
 H^0(K_v,A_{\eta,e})
       \lhook\joinrel\longrightarrow
 (O_\ell\otimes\mathbb Q_\ell/\mathbb Z_\ell(-e))^{G_{H_w}} .
 \tag{FL.13}
\]
The cyclotomic image on the right has index at most
\([H_w:\mathbb Q_\ell]\leq[H:\mathbb Q]=2D\).
The coefficient ring has trivial scalar Galois action,
and finite free extension by \(O_\ell\) commutes with
invariants.  Equation (FL.6) therefore bounds the
normalized order by
\[
 \frac1{b_\ell}\log\#H^0(K_v,A_{\eta,e})
                    \leq\log(8eD).
 \tag{FL.14}
\]
This group is zero unless \(\ell\leq8eD\).
There are \(D\) characters and at most two places
of \(K\) above any rational prime.
Using the elementary upper bound \(8eD\) for the number
of relevant primes, their total contribution is at most
\[
                         16eD^2\log(8eD).
 \tag{FL.15}
\]

\emph{4. The other ramified places.}
Fix \(v\in\mathcal V\) and \(\ell\ne\operatorname{char}k_v\).
The cyclotomic character is trivial on \(I_v\).
If \(\eta\) is unramified at \(v\), then
\(A_{\eta,e}^{I_v}=A_{\eta,e}\) and Frobenius acts by
\(\eta(g)(Nv)^{-e}\), where \(g\in G\) is any lift of
Frobenius in the decomposition group.
This scalar is not one: \(e\geq1\), \(Nv>1\), and
\(\eta(g)\) is a root of unity.
Multiplication by its nonzero difference from one is
surjective on the divisible module
\(\operatorname{Frac}(O_\ell)/O_\ell\).
Hence in this case
\[
                         H^1_{\rm unr}(K_v,A_{\eta,e})=0.
 \tag{FL.16}
\]
The invariants are finite, and their normalized order is
\[
 \frac1{b_\ell}\log\#H^0(K_v,A_{\eta,e})
     =v_\ell\bigl((Nv)^e-\eta(g)\bigr)\log\ell .
 \tag{FL.17}
\]
The valuation is normalized by \(v_\ell(\ell)=1\).
Indeed for any nonzero integral \(\alpha\in O_\ell\),
the kernel of multiplication by \(\alpha\) on
\(\operatorname{Frac}(O_\ell)/O_\ell\) has order
\(\ell^{b_\ell v_\ell(\alpha)}\).
The factor \((Nv)^e\) is an \(\ell\)-adic unit here.

Let \(m\) be the order of the chosen \(g\) in \(G\).
Every \(m\)-th root of unity occurs \(D/m\) times
among the values \(\eta(g)\), whence
\[
             \prod_{\eta\in\widehat G}(X-\eta(g))
                            =(X^m-1)^{D/m}.
 \tag{FL.18}
\]
Although only the unramified characters were used in
(FL.17), their valuation sum is at most that over all
characters, because all factors are integral.
Sum (FL.17) first over characters and then over
\(\ell\ne\operatorname{char}k_v\).
Equation (FL.18) bounds the result by
\[
           \frac Dm\log\bigl((Nv)^{em}-1\bigr)
                            \leq De\log Nv .
 \tag{FL.19}
\]
The lift \(g\) was fixed at this place; it need not be
canonical.  Changing it by inertia changes none of
the unramified-character values needed for the bound.

If \(\eta\) is ramified at \(v\), choose inertia with
nontrivial character value \(\zeta\).
Then \(A_{\eta,e}^{I_v}\) is contained in the kernel
of \(\zeta-1\).
Unless \(\zeta\) has pure \(\ell\)-power order, this
difference is an \(\ell\)-adic unit, so the invariants
vanish.  If the inertia character has any nontrivial
prime-to-\(\ell\) part, choose inertia detecting that part
to get the same conclusion.
The only possible nonzero cases therefore have
\(\ell\mid D\), namely \(\ell=2,3\).
For a nontrivial \(\ell\)-power root of unity,
\[
                 v_\ell(\zeta-1)\leq\frac1{\ell-1}.
 \tag{FL.20}
\]
To see the exact value, a primitive root of order
\(\ell^j\) satisfies
\(\prod_{(a,\ell)=1}(1-\zeta^a)=\ell\);
all its \(\varphi(\ell^j)\) factors have the same
valuation.  Thus that valuation is
\(1/\varphi(\ell^j)\leq1/(\ell-1)\).
Both \(H^0(K_v,A_{\eta,e})\) and the quotient in
(FL.9) have order at most \(\#A_{\eta,e}^{I_v}\).
Their total normalized contribution, summed over
characters, both bad primes and all these places, is
at most \(2Dr c_0\).
This proves local finiteness and, with
(FL.12), (FL.15), and (FL.19), proves (FL.4).

The sums in question have finite support.
The coefficient-prime local terms have the finite
cutoff in (FL.14); the unramified terms at a fixed
ramified place are supported on prime divisors of the
nonzero integer in (FL.18); and the ramified-character
terms involve only \(2,3\).
AC-K25 already gives finite prime support for the
strict groups at each fixed \(H,e\).
Outside the union of these finite sets, (FL.11)--(FL.12)
force the global \(H^2\) groups to vanish.
No exchange of nonconvergent prime sums has occurred.

\emph{5. The actual source weight window.}
Since \(K\) is quadratic, the sum of \(\log Nv\)
over all its places above a rational prime \(p\) is
at most \(2\log p\).
Only \(2,3,7,q,d,269\) can ramify in this family, so
\[
 \mathcal R\leq
 2(\log2+\log3+\log7+\log q+\log d+\log269).
 \tag{FL.21}
\]
The window bounds used in (FG.11) give
\(\mathcal R=O_\varepsilon(x^{\varepsilon/5}(\log x)^2)\)
and \(D=O_\varepsilon(x^{2+8\varepsilon})\).
Summing (FL.4) over \(e<t_x\) is at most a constant times
\[
 D^2t_x^2\log(8t_xD)+D\mathcal R t_x^2+Drt_x.
 \tag{FL.22}
\]
Its first term has order
\(O_\varepsilon(x^{806+16\varepsilon}\log x)\);
the second has order
\(O_\varepsilon(x^{804+41\varepsilon/5}(\log x)^2)\),
and the third is smaller.
This proves (FL.5), since every fixed polynomial in
\(x,\log x\) is little-o of \(\exp(x^\varepsilon/2)\)
for fixed \(\varepsilon>0\).
\end{proof}

\begin{proposition}[AC-K26-torsion: the all-prime degree-one torsion cost]
\label{AC-K26-torsion}
With the same coefficients and \(e\geq1\),
\[
 H^1(G_{K,\Sigma_\ell},T_{\eta,e})_{\rm tors}
       \simeq
 H^0\!\left(G_{K,\Sigma_\ell},
       (\operatorname{Frac}(O_\ell)/O_\ell)
                         (\eta^{-1})(e+1)\right).
 \tag{LT.1}
\]
This is a finite group.
Its normalized order summed over all primes and
characters satisfies
\[
 \sum_\ell\frac1{b_\ell}\sum_{\eta\in\widehat G}
        \log\#H^1(G_{K,\Sigma_\ell},T_{\eta,e})_{\rm tors}
       \leq8(e+1)D^2\log(8(e+1)D).
 \tag{LT.2}
\]
The sum is supported on \(\ell\leq8(e+1)D\).
Over \(1\leq e<t_x\), this again contributes
\(O_\varepsilon(x^{806+16\varepsilon}(\log x)^2)\),
and hence an error smaller than every fixed positive
fraction of the source exponential budget for large \(x\).
\end{proposition}
\begin{proof}
Put \(V_{\eta,e}=T_{\eta,e}\otimes_{\mathbb Z_\ell}\mathbb Q_\ell\).
Its global invariants are zero: restriction to \(G_H\)
removes \(\eta\), and the positive cyclotomic power
\(e+1\) is nontrivial on an open cyclotomic image.
The continuous cohomology sequence of
\[
              0\longrightarrow T_{\eta,e}
               \longrightarrow V_{\eta,e}
               \longrightarrow V_{\eta,e}/T_{\eta,e}
               \longrightarrow0
 \tag{LT.3}
\]
therefore injects \(H^0(V_{\eta,e}/T_{\eta,e})\)
into \(H^1(T_{\eta,e})\).
Its image is the kernel of the map to \(H^1(V_{\eta,e})\).
That kernel equals the torsion subgroup: its image from
the invariant quotient is torsion, while every torsion
class maps to zero in the \(\mathbb Q_\ell\)-vector space
\(H^1(V_{\eta,e})\).
This proves (LT.1), including the direction of the map.

Restriction of invariants to \(G_H\) is injective.
The latter group acts through a cyclotomic subgroup
of index at most \(2D\).
Apply (FL.6) with \(w=e+1\), and extend scalars by the
finite free \(O_\ell\).
The normalized logarithm is at most
\(\log(8(e+1)D)\), and it is zero unless
\(\ell\leq8(e+1)D\).
There are \(D\) characters and at most \(8(e+1)D\)
possible coefficient primes.
This gives (LT.2).
Its finite-weight sum has the same leading bound
\(D^2t_x^2\log(8t_xD)\) as (FL.22), which proves
the last assertion.
\end{proof}

\paragraph{What remains arithmetic.}
Equations (FL.4)--(FL.5) and (LT.2) remove the local
invariant terms, the inertia-versus-full-localization
difference and the degree-one torsion from the prescribed
budget at a proved polynomial cost.
They do not bound the determinant index or the free
degree-one regulator that controls the remaining global
degree-two sizes.
An estimate for those actual fixed-base character
quantities, with the common growing-host and weight
quantifiers, is still required.
The duality sources used above are available at
\url{https://math.ucla.edu/~sharifi/notes/iwasawa-app-a.html};
their displayed compact extensions, not a new proof of
global duality, are the background input to this section.
